\section{Discussion}

The main claim of this work is conceptual before it is algorithmic. Fourier-operator metasurface design is often narrated as a sequence of special cases: one paper for differentiation, another for edge detection, another for polarization multiplexing, another for spectral shaping. That style of narration tracks applications, but it can obscure the deeper common structure of the problem. Once the target is written as a coupled response field over angle, wavelength, and polarization, these tasks appear as different cuts through the same manifold. The advantage of the present formulation is therefore not merely that one generator can be reused across cases. The deeper advantage is that the optical specification and the inverse-design procedure are written in the same physical language.

This viewpoint also clarifies why engineering realism matters. Freeform metasurface inverse design is constrained not only by model capacity, but by electromagnetic cost, sparse feasible sets, and the need to retain fabrication-relevant structure. The repository reflects these realities directly: RCWA labeling is expensive, the forward surrogate is introduced to reduce but not eliminate that burden, the diffusion model is trained with an explicit binarization-related term, and final decisions are postponed until physical reranking. These are not implementation details to be hidden behind polished diagrams. They are the practical facts that make a physics-consistent generative treatment necessary in the first place.

Several limitations remain and should be stated plainly. First, the present response representation focuses on magnitude responses in two polarization channels. This is sufficient for the operator families considered here, but some metasurface functions require full complex-field control, phase-sensitive objectives, or a richer polarization basis. Second, the response field is discretized on a finite wavelength--angle grid. Sharp resonant features, narrow angular lobes, and subtle off-grid behavior may therefore be underresolved. Third, the geometry class is fabrication-aware only to a limited degree: binarization and symmetry are encoded, but a complete manufacturing model would also need to consider line-edge roughness, minimum-gap rules, etch bias, and tolerance under parameter drift.

These limitations do not weaken the central thesis so much as delimit its current scope. The response-space formulation is broad enough to absorb richer observables, denser sampling, and stronger fabrication constraints without changing its basic logic. Full complex responses could be introduced by expanding the condition channels. Robustness-aware objectives could be implemented by scoring candidates over perturbed material or geometric ensembles rather than over a single nominal field. Tolerance-aware generation could be incorporated by redefining the feasible set to include manufacturability margins alongside nominal optical fidelity. In each case the same principle would remain intact: the inverse problem should be stated and solved in the space of physically meaningful response fields, and final claims of performance should be made only after the generated candidates are returned to the physical solver.
